\documentclass[12pt,a4paper]{amsart}

\usepackage{amsmath,amssymb,amsthm}
\usepackage{mathtools}
\usepackage{booktabs}
\usepackage{array}
\usepackage[ngerman]{babel}
\usepackage{hyperref}
\usepackage{geometry}
\geometry{margin=2.5cm}

% -------------------------------------------------------
% Theorem-Umgebungen (deutsch)
% -------------------------------------------------------
\newtheorem{theorem}{Satz}[section]
\newtheorem{lemma}[theorem]{Lemma}
\newtheorem{korollar}[theorem]{Korollar}
\newtheorem{proposition}[theorem]{Proposition}
\theoremstyle{definition}
\newtheorem{definition}[theorem]{Definition}
\newtheorem{bemerkung}[theorem]{Bemerkung}

% -------------------------------------------------------
% Eigene Befehle
% -------------------------------------------------------
\newcommand{\N}{\mathbb{N}}
\newcommand{\Z}{\mathbb{Z}}
\newcommand{\Q}{\mathbb{Q}}
\newcommand{\R}{\mathbb{R}}
\newcommand{\C}{\mathbb{C}}
\newcommand{\F}{\mathbb{F}}
\DeclareMathOperator{\li}{li}
\newcommand{\eexp}[1]{e\!\left(#1\right)}
\newcommand{\Ptwo}{\mathcal{P}_2}
\newcommand{\Pk}{\mathcal{P}_k}
\DeclareMathOperator{\ord}{ord}
\newcommand{\OOmega}{\Omega}

% -------------------------------------------------------
\begin{document}
% -------------------------------------------------------

\title{Chens Satz: Jede hinreichend große gerade Zahl\\
       ist Summe einer Primzahl und einer Fastprimzahl}

\author{Michael Fuhrmann}
\address{Specialist Mathematics Project, Build 56}
\email{specialist-maths@localhost}
\date{2026}

\begin{abstract}
Wir stellen Chens berühmten Satz (1966, veröffentlicht 1973) vor:
Jede hinreichend große gerade ganze Zahl $N$ lässt sich schreiben als
$N = p + q$,
wobei $p$ eine Primzahl und $q \in \Ptwo$ ist (d.h., $q$ ist entweder
prim oder ein Produkt von genau zwei Primzahlen, mit Vielfachheit gezählt).
Dies ist das der Goldbachschen Vermutung am nächsten bekannte Ergebnis.
Wir geben eine vollständige Beweisskizze unter Verwendung des Selberg-Siebs
in Kombination mit dem Buchstab-Umschaltprinzip und Rosser--Iwaniec-Gewichten
an und beweisen die untere Schranke
\[
  \#\bigl\{p \leq N : N - p \in \Ptwo\bigr\}
  \;\geq\;
  \frac{c\, N}{(\log N)^2}\prod_{p \mid N}\frac{p-1}{p-2}
\]
für eine absolute Konstante $c > 0$.  Wir erläutern das ``Paritätsproblem'',
das eine Ausdehnung der Methode auf einen Beweis der Goldbachschen Vermutung
verhindert.
\end{abstract}

\maketitle
\tableofcontents

% ================================================================
\section{Einleitung: Von Goldbach zu Chen}
% ================================================================

Die Goldbachsche Vermutung (1742) behauptet, dass jede gerade ganze Zahl
$N \geq 4$ Summe zweier Primzahlen ist.  Trotz ihrer Stellung als eines der
ältesten ungelösten Probleme der Mathematik hat sie alle Beweisversuche
widerstanden.

Die Siebbmethoden des 20.~Jahrhunderts lieferten eine Reihe sukzessiv
stärkerer Teilresultate, gemessen am \emph{Typ} $(k, \ell)$: Jede hinreichend
große gerade Zahl $N$ ist Summe einer Zahl mit höchstens $k$ Primfaktoren
und einer Zahl mit höchstens $\ell$ Primfaktoren.

\begin{center}
\begin{tabular}{@{}lll@{}}
\toprule
Jahr & Autor & Ergebnis \\ \midrule
1919 & Brun        & $(9,9)$: Jedes große gerade $N = a + b$, $\OOmega(a),\OOmega(b) \leq 9$ \\
1937 & Buchstab    & $(5,5)$ \\
1940 & Buchstab    & $(4,4)$ \\
1956 & Wang Yuan   & $(3,4)$ \\
1957 & Wang Yuan   & $(2,3)$ \\
1962 & Barban, Pan & $(1,5)$ \\
1965 & Buchstab    & $(1,3)$ \\
1966 & Chen Jingrun & $(1,2)$: jedes große gerade $N = p + q$, $\OOmega(q) \leq 2$ \\
\bottomrule
\end{tabular}
\end{center}

Die Zielsequenz ist $(1,\infty) \to (1,3) \to (1,2) \to (1,1)$, wobei
der letzte Schritt $(1,1)$ die Goldbachsche Vermutung ist.  Chens Satz
erreicht $(1,2)$ --- das Bestmögliche, das mit der aktuellen Siebtheorie
erreichbar ist, aufgrund des ``Paritätsproblems'' (vgl. Abschnitt~\ref{sec:paritaet}).

Chen kündigte sein Ergebnis 1966 an; der vollständige Beweis erschien 1973
\cite{Chen1973}.

% ================================================================
\section{Fastprimzahlen und die $\Ptwo$-Notation}
% ================================================================

\begin{definition}[Fastprimzahl der Ordnung $k$]
  Für $k \geq 1$ heißt eine positive ganze Zahl $n$ eine
  \emph{$\mathcal{P}_k$-Zahl} (oder \emph{Fastprimzahl der Ordnung $k$}), falls
  \[
    \OOmega(n) \;\leq\; k,
  \]
  wobei $\OOmega(n) = \sum_{p^a \| n} a$ die Primfaktoren mit Vielfachheit
  zählt.  Wir schreiben $\Ptwo$ für die Menge der $\mathcal{P}_2$-Zahlen.
\end{definition}

Es gilt also $\Ptwo = \{\text{Primzahlen}\} \cup \{pq : p, q \text{ Primzahlen}\}$.
Zum Beispiel enthält $\Ptwo$ die Zahlen $2, 3, 4, 5, 6, 7, 9, 10, 11, 13,
14, 15, \ldots$ (alle Primzahlen und Semiprimes).

\begin{lemma}[Zählen von Fastprimzahlen]\label{lem:zaehlen_p2}
  Die Anzahl $\omega_2(x) := \#\{n \leq x : n \in \Ptwo\}$ erfüllt
  \[
    \omega_2(x) \;\sim\; \frac{x\log\log x}{2\log x}
    \quad (x \to \infty).
  \]
\end{lemma}

\begin{proof}
  Durch partielle Summation und Mertens' Satz.  Die Primzahlen bis $x$
  liefern den Beitrag $\sim x/\log x$.  Die Semiprimes $pq \leq x$ mit
  $p \leq q$ tragen, durch zweifache Anwendung des Primzahlsatzes, bei:
  \[
    \sum_{p \leq \sqrt{x}}\pi(x/p)
    \sim \sum_{p \leq \sqrt{x}}\frac{x/p}{\log(x/p)}
    \sim \frac{x}{\log x}\sum_{p \leq \sqrt{x}}\frac{1}{p}
    \sim \frac{x\log\log x}{2\log x},
  \]
  wobei der letzte Schritt Mertens' Satz $\sum_{p \leq y}1/p \sim \log\log y$
  verwendet.  Der Beitrag der Primzahlen $x/\log x = o(x\log\log x/\log x)$
  ist vernachlässigbar.
\end{proof}

% ================================================================
\section{Der Siebansatz für Chens Satz}
% ================================================================

Sei $N$ eine hinreichend große gerade ganze Zahl.  Definiere die Menge der
verschobenen Primzahlen:
\[
  \mathcal{A} \;:=\; \{N - p : p \leq N,\; p \text{ prim}\},
\]
so gilt $|\mathcal{A}| \approx N/\log N$ nach dem Primzahlsatz.
Wir möchten zeigen, dass $\mathcal{A}$ Elemente in $\Ptwo$ enthält.

\begin{definition}[Siebfunktion]
  Für eine Menge $\mathcal{A} \subseteq \N$, eine Primzahlmenge $\mathcal{P}$
  und einen Parameter $z > 0$ definiere
  \[
    S(\mathcal{A}, \mathcal{P}, z)
    \;:=\;
    \#\bigl\{a \in \mathcal{A} : \ggT(a, P(z)) = 1\bigr\},
    \quad P(z) = \prod_{\substack{p \in \mathcal{P} \\ p < z}} p.
  \]
\end{definition}

Für $z = N^{1/2}$ und $\mathcal{P}$ gleich aller Primzahlen zählt
$S(\mathcal{A}, \mathcal{P}, N^{1/2})$ die Elemente von $\mathcal{A}$
ohne Primteiler unterhalb von $N^{1/2}$ --- das sind genau die Primzahlen
in $\mathcal{A}$.

Für einen Teiler $d \mid P(z)$ bezeichne $\mathcal{A}_d := \{a \in \mathcal{A} : d \mid a\}$.
% BUG-P16-005 behoben (DE): ω(p) war als p/(p-1) angegeben, was einer anderen
% Normierungskonvention entspricht. In der Darstellung |A_d|=(ω(d)/d)|A|+r_d
% gilt ω(p)=1 für p∤N (jede Restklasse mod p trifft ~1/p der Primes,
% also |A_p| ≈ |A|/p, d.h. ω(p)/p = 1/p → ω(p)=1). Vereinheitlicht auf ω(p)=1.
Die multiplikative Funktion $\omega$ für unser $\mathcal{A}$ erfüllt für
$p \nmid N$: $\omega(p) = 1$ (da ungefähr $1/p$ der Primzahlen in jede
Restklasse $\bmod p$ fällt, also $|\mathcal{A}_p| \approx |\mathcal{A}|/p$,
d.h. $\omega(p)/p \approx 1/p$, also $\omega(p) = 1$).

% ================================================================
\section{Die Buchstab-Identität und das Umschaltprinzip}
% ================================================================

Die Buchstab-Identität ist das induktive Rückgrat des Siebs:

\begin{proposition}[Buchstab-Identität]\label{prop:buchstab}
  Für $w < z$ gilt:
  \[
    S(\mathcal{A}, \mathcal{P}, z)
    \;=\;
    S(\mathcal{A}, \mathcal{P}, w)
    \;-\;
    \sum_{\substack{p \in \mathcal{P} \\ w \leq p < z}}
    S(\mathcal{A}_p, \mathcal{P}, p).
  \]
\end{proposition}

\begin{proof}
  Ein Element $a \in \mathcal{A}$ wird von $S(\mathcal{A},\mathcal{P},w)$,
  aber nicht von $S(\mathcal{A},\mathcal{P},z)$ gezählt, genau dann, wenn
  $a$ einen Primteiler $p \in \mathcal{P}$ mit $w \leq p < z$ besitzt.
  Sei $p_0$ der \emph{kleinste} solche Primteiler.  Dann gilt
  $a \in \mathcal{A}_{p_0}$, und $a/p_0$ hat keinen Primteiler in
  $\mathcal{P}$ unterhalb von $p_0$, also trägt $a/p_0$ zu
  $S(\mathcal{A}_{p_0}, \mathcal{P}, p_0)$ bei.  Summation über alle
  solchen $p_0 = p$ liefert die Formel.  \qedhere
\end{proof}

Das \textbf{Umschaltprinzip} (Chens Schlüsselidee) wendet die
Buchstab-Identität mit $z = N^{1/2}$ und $w = N^{1/3}$ an:
\begin{align}
  S(\mathcal{A}, \mathcal{P}, N^{1/2})
  &= S(\mathcal{A}, \mathcal{P}, N^{1/3})
  - \sum_{N^{1/3} \leq p < N^{1/2}} S(\mathcal{A}_p, \mathcal{P}, p).
  \label{eq:buchstab_split}
\end{align}

Der zweite Term zählt Elemente $a \in \mathcal{A}$, deren kleinster
Primteiler $p \in [N^{1/3}, N^{1/2})$ ist.  Da $a = N - q_1$ mit $q_1$
prim, und $a$ hat einen Primteiler $p \geq N^{1/3}$, gilt $a = p \cdot m$
mit $m < N^{2/3}$.  Ist $m$ prim, so gilt $\OOmega(a) = 2$ und
$a \in \Ptwo$.

Die Elemente, die vom zweiten Term in \eqref{eq:buchstab_split} gezählt
werden und \emph{nicht} in $\Ptwo$ liegen, haben $\OOmega(m) \geq 2$,
also $a = pm'p'$ mit $p' \leq m^{1/2} \leq N^{1/3}$.  Man kann Buchstab
erneut anwenden, um diesen Fehler abzuschätzen.

% ================================================================
\section{Die Hauptabschätzung: Chens Satz}
% ================================================================

\begin{theorem}[Chen, 1973]\label{thm:chen}
  Für jede hinreichend große gerade ganze Zahl $N$ gibt es eine Primzahl
  $p$ und eine ganze Zahl $q \in \Ptwo$ mit $N = p + q$.

  Quantitativ: Mit der singulären Reihe
  \[
    \mathfrak{S}(N)
    \;:=\;
    2\prod_{\substack{p \geq 3 \\ p \mid N}}
    \frac{p-1}{p-2}
    \cdot
    \prod_{\substack{p \geq 3 \\ p \nmid N}}
    \left(1 - \frac{1}{(p-1)^2}\right)
  \]
  (die $\mathfrak{S}(N) > 0$ für alle geraden $N$ erfüllt) gilt:
  \[
    \#\bigl\{p \leq N : N-p \in \Ptwo\bigr\}
    \;\geq\;
    \frac{0{,}67\,\mathfrak{S}(N)\, N}{(\log N)^2}
  \]
  für alle hinreichend großen geraden $N$.
\end{theorem}

\textit{Beweisskizze.}  Der Beweis kombiniert drei Bestandteile:

\textbf{Bestandteil 1: Selbergs Sieb-Oberschranke.}
Das Selberg-Sieb liefert eine Oberschranke für $S(\mathcal{A},\mathcal{P},z)$.
Mit der Siebbdimension $\kappa = 1$ (da $\omega(p) \approx 1$ im Mittel),
dem Siebniveau $D = N^{1/2}/(\log N)^B$ und dem Siebgrenzwert
$\beta = (\log D)/(\log z)$ erhält man:
\[
  S(\mathcal{A}, \mathcal{P}, N^{1/2})
  \;\leq\;
  (F(\beta) + \varepsilon)\frac{W(z) N}{\phi(N/2)\log D}
  + O\!\left(\frac{N}{(\log N)^A}\right),
\]
wobei $W(z) = \prod_{p < z}(1 - \omega(p)/p)$ das Siebprodukt ist und
$F$ die obere Siebfunktion mit $F(s) \leq 2e^\gamma/s$ für $s > 1$.

\textbf{Bestandteil 2: Rosser--Iwaniec-Unterschranke.}
Für die untere Siebabschätzung verwendet man Rossers kombinatorisches Sieb
(mit Iwaniec-Verbesserungen), das eine \emph{untere} Schranke liefert:
\[
  S(\mathcal{A}, \mathcal{P}, N^{1/3})
  \;\geq\;
  (f(\beta) - \varepsilon)\frac{W(z)N}{\phi(N/2)\log D}
  - O\!\left(\frac{N}{(\log N)^A}\right),
\]
wobei $f$ die untere Siebfunktion mit $f(s) \geq 0$ für $s > 2$ ist.
Speziell gilt $f(s) = 2e^\gamma \log(s-1)/s$ für $s \in (2, 4]$.

\textbf{Bestandteil 3: Das Buchstab-Umschalten.}
Unter Verwendung von \eqref{eq:buchstab_split} und Anwendung des Siebs
auf die inneren Summen ergibt sich:
\begin{equation}\label{eq:chen_schluessel}
  \#\{p \leq N : N-p \in \Ptwo\}
  \;\geq\;
  S(\mathcal{A}, \mathcal{P}, N^{1/3})
  - \tfrac{1}{2}\sum_{N^{1/3} \leq p < N^{1/2}} S(\mathcal{A}_p, \mathcal{P}, p) + \cdots
\end{equation}
% BUG-P16-006 behoben (DE): Faktor 1/2 wurde falsch als "Doppelzählung der
% Primteiler" erklärt. Korrektur: Der Faktor 1/2 kommt aus dem Umschaltprinzip
% (Switching Principle) — beim Tauschen der Rollen von p und q in der Summe
% wird jedes Paar (p, N-p) auf jeder Seite einmal gezählt; der Durchschnitt
% der beiden Abschätzungen liefert den Faktor 1/2.
Der Faktor $\tfrac{1}{2}$ entsteht durch das Umschaltprinzip: Wenn man die
Rollen von $p$ und $q$ in der Summe tauscht, wird jedes Paar $(p, N-p)$
mit $N-p \in \Ptwo$ auf jeder Seite einmal gezählt.  Genauer: Die
Buchstab-Zerlegung \eqref{eq:buchstab_split} wird zweimal angewendet (einmal
mit vertauschten Summanden), und das Mitteln der beiden Abschätzungen liefert
den Faktor $\tfrac{1}{2}$; vgl.\ Iwaniec--Kowalski \cite{IwaniecKowalski2004},
Kapitel~11.

Nach sorgfältiger Abschätzung aller Fehlerterme und partieller Summation
erhält man die untere Schranke mit der Konstante $c = 0{,}67$.  \qed

\begin{bemerkung}
  Die Konstante $0{,}67$ wurde von Chen \cite{Chen1973} selbst
  estabelliert und später von anderen Autoren verfeinert.  Der derzeit
  beste Wert beträgt ungefähr $0{,}748$ (Pan Chengbiao, 1992).
\end{bemerkung}

\begin{korollar}[Goldbach--Chen für große $N$]\label{kor:goldbach_chen}
  Für jedes gerade $N \geq N_0$ (für ein effektives $N_0$) hat die Gleichung
  $N = p + m$ mit $p$ prim und $\OOmega(m) \leq 2$ mindestens
  $\gg N/(\log N)^2$ Lösungen.  Insbesondere hat $N = p + m$ mindestens
  eine Lösung mit $p$ prim und $m \in \Ptwo$.
\end{korollar}

% ================================================================
\section{Numerische Verifikation}
% ================================================================

Wir verifizieren Chens Satz für kleine gerade Werte.  Für jedes gerade $N$
listen wir alle Darstellungen $N = p + q$ mit $p$ prim und $q \in \Ptwo
= \{\text{prim oder Semiprim}\}$ auf.  Zum Vergleich führen wir auch
Goldbach-Zerlegungen auf (wo beide Summanden prim sind).

% BUG-P16-001/002/003 (DE): Tabelle war bereits korrekt (keine 1∉P₂-Einträge).
% Bestätigt: Alle Einträge p+q mit q∈P₂ sind korrekt; 5+1, 7+1, 11+1 sind nicht aufgeführt.
\begin{center}
\begin{tabular}{@{}r l l@{}}
\toprule
$N$ & Chen-Zerlegungen $N = p + q$, $q \in \Ptwo$ & Goldbach ($q$ prim) \\ \midrule
4   & $4 = 2 + 2$                                   & $2 + 2$ \\
6   & $6 = 3 + 3$                                   & $3 + 3$ \\
8   & $8 = 3 + 5$                                   & $3 + 5$ \\
10  & $10 = 3 + 7 = 5 + 5 = 7 + 3$                & $3+7$, $5+5$ \\
12  & $12 = 5 + 7 = 7 + 5$                         & $5+7$, $7+5$ \\
14  & $14 = 3 + 11 = 7 + 7 = 11 + 3$              & $3+11$, $7+7$, $11+3$ \\
16  & $16 = 3 + 13 = 5 + 11 = 13 + 3$             & $3+13$, $5+11$, $13+3$ \\
18  & $18 = 5 + 13 = 7 + 11 = 11 + 7 = 13 + 5$   & alle aufgelisteten \\
20  & $20 = 3 + 17 = 7 + 13 = 13 + 7 = 17 + 3$   & alle aufgelisteten \\
\bottomrule
\end{tabular}
\end{center}

Für größere $N$ tauchen Chen-Zerlegungen mit echten Semiprimes auf:
\begin{align*}
  100 &= 5 + 95\;(95 = 5 \cdot 19 \in \Ptwo\;{\checkmark}) \\
      &= 7 + 93\;(93 = 3 \cdot 31 \in \Ptwo\;{\checkmark}) \\
      &= 11 + 89\;(89 \text{ prim}\;{\checkmark}) \\
      &= 3 + 97\;(97 \text{ prim}\;{\checkmark}).
\end{align*}

Die asymptotische Formel prognostiziert ungefähr
$0{,}67 \cdot \mathfrak{S}(100) \cdot 100/(\log 100)^2 \approx 4{,}6$
Darstellungen, was mit der beobachteten Anzahl übereinstimmt.

% ================================================================
\section{Die Lücke zu Goldbach: Das Paritätsproblem}\label{sec:paritaet}
% ================================================================

Warum kann Chens Methode nicht zur Goldbachschen Vermutung ausgedehnt werden?
Das Hindernis ist Selbergs \emph{Paritätsproblem}.

\begin{bemerkung}[Das Paritätsproblem, Selberg 1949]
  Betrachte das Sieben der Menge $\mathcal{A} = \{N - p : p \leq N\}$,
  um Elemente mit $\OOmega = 1$ (Primzahlen) vs. $\OOmega = 2$ (Semiprimes)
  zu finden.  Diese beiden Mengen haben dieselbe ``Sieb-Signatur'': Für
  jeden quadratfreien Modul $d$ ist die Dichte von $\mathcal{A}_d$ in
  $\mathcal{A}$ dieselbe, egal ob wir $\OOmega = 1$ oder $\OOmega = 2$
  fordern.  Das Sieb kann sie nicht unterscheiden.

  Formal: Die Möbius-Funktion $\mu(n)$ hat die Paritätseigenschaft
  $\mu(n) = (-1)^{\OOmega(n)}$ für quadratfreie $n$.  Ein reines Sieb
  tastet $\mathcal{A}$ nur über multiplikative Funktionen von
  $\ggT(a, P(z))$ ab, und $\OOmega(a)$ modulo~$2$ ist für solche
  Abtastungen unsichtbar.

  Die charakteristische Funktion von ``$\OOmega(n)$ ist gerade'' und die
  von ``$\OOmega(n)$ ist ungerade'' haben dieselbe multiplikative Struktur
  modulo Quadratfreier; daher weist ein reines kombinatorisches Sieb beiden
  dieselbe untere Schranke (null) oder Oberschranke zu.
\end{bemerkung}

Was wäre nötig, um von Chens $(1,2)$ zu Goldbachs $(1,1)$ zu gelangen?
Man benötigt analytische Zusatzinformation jenseits der Siebkombinatorik,
z.B.:
\begin{itemize}
  \item Information über die Verteilung von Primzahlen in kurzen Intervallen
    (dies würde die GRH oder ihre Konsequenzen erfordern),
  \item Ein effektives Mittel zur Unterscheidung der $\OOmega$-Parität im
    Sieb, das die aktuelle Technologie nicht bereitstellen kann,
  \item Einen grundlegend neuen Ansatz (etwa über Exponentialsummen,
    die Kreismethode oder Methoden aus Zahlkörpern).
\end{itemize}

Die Goldbachsche Vermutung für Primzahlen im Intervall $[N/3, N/2]$ ist
offen; ein Nachweis, dass $\{N - p : p \in [N/3, N/2]\}$ eine Primzahl
enthält, würde Goldbach liefern.  Dies ist das ``kurze-Intervall-Goldbach-
Problem'', das ebenfalls offen ist.

% ================================================================
\section{Schlussfolgerung}
% ================================================================

Chens Satz stellt den Höhepunkt der Siebtheorie des 20.~Jahrhunderts
angewandt auf das Goldbach-Problem dar.  Sein Beweis vereint:
\begin{enumerate}
  \item \textbf{Das Selberg-Sieb} für Oberschranken, das das ``Verteilungsniveau''
    von Primzahlen in arithmetischen Progressionen bereitstellt,
  \item \textbf{Das Rosser--Iwaniec-Linearsieb} für Unterschranken,
    das die eindimensionale Struktur des Problems ausnutzt,
  \item \textbf{Die Buchstab-Identität} für das Umschaltprinzip, das das
    Problem vom Nachweis von Primzahlen auf den Nachweis von $\Ptwo$-Zahlen
    reduziert,
  \item \textbf{Die singuläre Reihe} $\mathfrak{S}(N)$, die die arithmetische
    Struktur des geraden $N$ kodiert und stets positiv ist.
\end{enumerate}

Das Ergebnis $N = p + q$ mit $\OOmega(q) \leq 2$ ist innerhalb des Rahmens
der reinen Siebmethoden das Bestmögliche, bedingt durch das Paritätshindernis.
Dieses Hindernis zu überwinden und die Goldbachsche Vermutung zu beweisen
bleibt eines der berühmtesten offenen Probleme der Mathematik.

\begin{thebibliography}{10}

\bibitem{Chen1973}
J.~Chen.
\newblock On the representation of a large even integer as the sum of a prime
  and the product of at most two primes.
\newblock \emph{Sci.\ Sinica}, 16:157--176, 1973.

\bibitem{MV1973}
H.~L.~Montgomery und R.~C.~Vaughan.
\newblock The large sieve.
\newblock \emph{Mathematika}, 20:119--134, 1973.

\bibitem{IwaniecKowalski2004}
H.~Iwaniec und E.~Kowalski.
\newblock \emph{Analytic Number Theory}.
\newblock American Mathematical Society, Providence, RI, 2004.

\bibitem{Bombieri1965}
E.~Bombieri.
\newblock On the large sieve.
\newblock \emph{Mathematika}, 12:201--225, 1965.

\bibitem{CojocMurty2005}
A.~C.~Cojocaru und M.~R.~Murty.
\newblock \emph{An Introduction to Sieve Methods and Their Applications}.
\newblock Cambridge University Press, Cambridge, 2005.

\bibitem{Halberstam1974}
H.~Halberstam und H.-E.~Richert.
\newblock \emph{Sieve Methods}.
\newblock Academic Press, London, 1974.

\bibitem{Goldbach1742}
C.~Goldbach.
\newblock Brief an L.~Euler, 7.~Juni 1742.
\newblock In: \emph{Briefwechsel von Euler}, 1742.

\bibitem{Brun1919}
V.~Brun.
\newblock Le crible d'Eratosth\`{e}ne et le th\'{e}or\`{e}me de Goldbach.
\newblock \emph{C.\ R.\ Acad.\ Sci.\ Paris}, 168:544--546, 1919.

\bibitem{Pan1992}
C.~Pan.
\newblock On the representation of a large even integer.
\newblock \emph{Chin.\ Ann.\ Math.\ B}, 1992.

\end{thebibliography}

\end{document}
